\section{Timeline Semantics and Cursor-Based Time}\label{sec:temporal-model}

A central challenge in music programming is the representation of time.
While physical time is continuous, musical time is often perceived as a series of discrete units (beats, bars,
subdivisions).
Motivo uses a deterministic \textit{cursor model} to manage temporal placement, providing a balance
between sequential flow and absolute positioning.

\subsection{The Virtual Playhead}\label{detail-subsec:the-virtual-playhead}

The core of the temporal logic is a virtual ``playhead'' or \textit{cursor}.
As the compiler traverses the musical statements within a track or voice, it maintains an internal state representing
the current absolute beat position.
Every musical command interacts with this cursor in one of two ways:

\noindent \textbf{Relative Advancement.} Standard \texttt{play} and \texttt{rest} statements are relative to the current
cursor position.
When a note is played, it is placed at the current cursor value, and the cursor is subsequently advanced by the
duration of that note.
This model allows for natural, sequential composition that mimics the act of writing on a staff.

\noindent \textbf{Absolute Offsets.} The \texttt{from} keyword allows for non-destructive positioning.
A statement such as \texttt{play A4 from 16;} places a note at the 16th beat regardless of the current cursor
position.
Crucially, this operation does not advance the main cursor, allowing the composer to ``jump'' back and forth across
the timeline to layer multiple events without losing the primary sequential flow.

\subsection{Parallelism and Temporal Isolation}\label{detail-subsec:parallelism-and-temporal-isolation}

The cursor model is also applied to parallel structures.
When a \texttt{voice} is defined within a \texttt{track}, it inherits a copy of the parent track's current cursor
position.
However, advancements within the voice are isolated; they do not affect the cursor of the parent track or of sibling
voices.
This allows for the creation of independent rhythmic layers that are synchronized at their start point but proceed at
their own pace.

\subsection{Deterministic Time Resolution}\label{detail-subsec:deterministic-time-resolution}

By using floating-point arithmetic for the internal cursor, the language supports arbitrary rhythmic subdivisions
(e.g., triplets, dotted notes, or complex swing rhythms).
Unlike interpreted systems that may suffer from timing drift or ``jitter'' due to CPU scheduling, the compiled nature of
Motivo calculates every event before writing the output file, so event placement is resolved before playback.
Time is resolved entirely during the ``Lowering'' phase of the compiler, transforming the hierarchical code into a flat,
sorted timeline of events.
This methodology keeps musical placement deterministic within the current mapping.

\begin{figure}[htbp]
  \centering
  \begin{tikzpicture}[
      beat/.style={font=\scriptsize\sffamily, text=motivoInk},
      note/.style={rounded corners=2pt, minimum height=0.42cm, draw=white, very thick},
      arrow/.style={-Stealth, thick, draw=motivoInk!70}
    ]
    \draw[thick, draw=motivoInk!70] (0,0) -- (12,0);
    \foreach \x/\b in {0/0,2/1,4/2,6/3,8/4,10/5,12/6} {
      \draw[draw=motivoInk!45] (\x,0.12) -- (\x,-0.12);
      \node[beat, below=0.18cm] at (\x,0) {\b};
    }

    \node[note, fill=motivoBlue!70, minimum width=1.9cm] at (1,0.65) {\textcolor{white}{play C4}};
    \node[note, fill=motivoGreen!70, minimum width=1.9cm] at (3,0.65) {\textcolor{white}{play D4}};
    \node[note, fill=motivoAmber!85, minimum width=1.9cm] at (5,0.65) {\textcolor{white}{play E4}};
    \node[note, fill=motivoRose!80, minimum width=1.9cm] at (9,1.45) {\textcolor{white}{play G4 from 4}};

    \draw[arrow] (0,1.1) -- node[beat, above] {relative cursor advances} (6,1.1);
    \draw[arrow] (9,1.15) -- node[beat, right, align=left] {\texttt{from}\\explicit placement\\cursor
    unchanged} (8.1,0.88);
  \end{tikzpicture}
  \caption{Cursor-based temporal placement.
    Sequential statements advance the cursor, while \texttt{from} positions an event on the timeline
  explicitly.}\label{fig:cursor-model}
\end{figure}
